\begingroup
\setlength{\abovedisplayskip}{7pt}
\setlength{\belowdisplayskip}{7pt}
\setlength{\abovedisplayshortskip}{3pt}
\setlength{\belowdisplayshortskip}{4pt}

% Printed 52 (unfoliated); principal PDF 59; Google PDF 63.
\par\clearpage

\vspace*{.19\textheight}
\begin{center}
{\large IV.}\par\medskip
{\Large LATITUDE DETERMINATION\\FOR THE SEEBERG OBSERVATORY.}
\end{center}

At the same time as I made my observations in Göttingen and Altona,
Mr. \textsc{Hansen}, director of the Seeberg Observatory near Gotha, also
observed the same stars at my request with the two-foot \textsc{Ertel} meridian
circle there.
The resulting difference in latitude between that observatory and the one
in Göttingen is of still greater interest because the former is connected to the
Hanoverian system of triangles by several triangles measured under the direction
of Lieutenant General \textsc{von Müffling}.

The circle was reversed several times during the observations.
Independently of these reversals, the collimation error was determined every day,
usually twice a day, by setting the instrument on the nadir point. Mr. \textsc{Hansen}
had learned the procedure already mentioned above in practice at this observatory
in the autumn of 1826. The readings were made with microscopes instead of verniers.
The following table gives the principal results of these observations. The first
column gives the designation of the star, the second the position of the circle,
the third the number of observations, the fourth the zenith distance (positive
to the north) as I reduced it to the beginning of 1827, and the fifth the latitude
that follows from the declinations given on p.~9 above.

% Printed 53; principal PDF 60; Google PDF 64.
\par\clearpage
\noindent\makebox[12mm][l]{}\parbox[t]{\dimexpr\linewidth-24mm\relax}{\centering\scriptsize\bfseries DETERMINATION OF THE DIFFERENCE IN LATITUDE ETC.\quad IV. LATITUDE DETERMINATION FOR SEEBERG.}\makebox[12mm][r]{\fontsize{13}{16}\selectfont 53}\par\medskip

{\small
\setlength{\arraycolsep}{3.2pt}
\renewcommand{\arraystretch}{1.06}
\[
\begin{array}{@{}r|c|c|r@{\;}r@{\;}r|r||r|c|c|r@{\;}r@{\;}r|r@{}}
1&\text{East}&5&-1^\circ&1'&53''{,}14&50^\circ56'4''{,}76&
16&\text{East}&5&-1^\circ&42'&19''{,}87&50^\circ56'6''{,}55\\
&\text{West}&4&&&52{,}80&4{,}42&&\text{West}&5&&&17{,}05&3{,}73\\
2&\text{East}&5&+4&37&29{,}41&5{,}57&17&\text{East}&5&+0&38&48{,}22&5{,}50\\
&\text{West}&6&&&30{,}86&4{,}12&&\text{West}&5&&&48{,}79&4{,}93\\
3&\text{East}&5&-0&45&19{,}00&5{,}20&18&\text{East}&5&+1&38&59{,}63&6{,}24\\
&\text{West}&7&&&17{,}70&3{,}90&&\text{West}&6&&&61{,}51&4{,}36\\
4&\text{East}&6&+3&38&50{,}38&5{,}17&19&\text{West}&2&+3&41&30{,}38&4{,}49\\
&\text{West}&6&&&51{,}05&4{,}50&20&\text{East}&1&+3&33&50{,}19&4{,}11\\
5&\text{East}&6&+4&29&54{,}35&4{,}87&&\text{West}&1&&&47{,}87&6{,}43\\
&\text{West}&6&&&54{,}30&4{,}92&21&\text{East}&1&+1&58&31{,}74&5{,}35\\
6&\text{East}&6&-4&20&26{,}87&6{,}01&&\text{West}&1&&&32{,}90&4{,}19\\
&\text{West}&4&&&25{,}92&5{,}06&22&\text{East}&1&-4&39&61{,}37&6{,}07\\
7&\text{East}&6&-0&39&22{,}12&5{,}56&&\text{West}&1&&&58{,}90&3{,}60\\
&\text{West}&4&&&20{,}70&4{,}14&23&\text{East}&1&+5&24&21{,}09&5{,}91\\
8&\text{East}&5&+1&40&1{,}41&6{,}06&&\text{West}&1&&&22{,}88&4{,}12\\
&\text{West}&4&&&3{,}24&4{,}23&24&\text{West}&1&-0&33&26{,}54&5{,}95\\
9&\text{East}&5&+5&17&29{,}69&6{,}82&25&\text{West}&1&+8&5&40{,}64&6{,}50\\
&\text{West}&4&&&33{,}72&2{,}79&26&\text{West}&1&-0&17&51{,}91&5{,}03\\
10&\text{East}&5&+1&43&6{,}51&5{,}54&27&\text{West}&1&-4&35&46{,}53&3{,}57\\
&\text{West}&5&&&8{,}63&3{,}42&28&\text{West}&1&+2&44&9{,}48&3{,}93\\
11&\text{East}&5&+2&43&27{,}76&5{,}26&29&\text{West}&1&+1&31&5{,}78&5{,}28\\
&\text{West}&4&&&28{,}89&4{,}13&30&\text{West}&2&+4&40&0{,}22&4{,}51\\
12&\text{East}&5&+2&2&48{,}89&6{,}77&31&\text{West}&1&-4&57&58{,}28&4{,}13\\
&\text{West}&3&&&50{,}29&5{,}37&32&\text{West}&1&+2&18&59{,}76&3{,}88\\
13&\text{East}&5&-1&29&56{,}73&6{,}19&33&\text{West}&1&-0&39&51{,}61&5{,}28\\
&\text{West}&5&&&56{,}64&6{,}10&34&\text{West}&1&+6&9&34{,}85&2{,}94\\
14&\text{East}&5&-0&35&42{,}51&5{,}00&35&\text{West}&1&-3&59&15{,}56&4{,}09\\
&\text{West}&3&&&42{,}84&5{,}33&36&\text{West}&2&-4&6&42{,}83&4{,}84\\
15&\text{East}&5&-2&36&14{,}46&6{,}93&37&\text{West}&1&+6&0&36{,}97&5{,}38\\
&\text{West}&4&&&12{,}88&5{,}35&38&\text{West}&1&-1&53&16{,}75&3{,}84
\end{array}
\]
\par
}

These sixty results for the latitude are, of course, of unequal reliability.
To assign weights to them without arbitrariness, one would have to know the ratio
of the mean observational error proper to the mean graduation error. If this ratio
is \(1\) to \(\surd\theta\), then, disregarding the slight uncertainty still
attaching to the declinations,
\[
\frac n{1+n\theta}
\]
\noindent will be the weight of a determination based on \(n\) observations
relating to one and the same graduation mark.

If one simply takes \(n\) in place of this weight, the mean of the 206
observations is
\[
=50^\circ56'5''{,}16.
\]

% Printed 54; principal PDF 61; Google PDF 65.
\par\clearpage
\noindent\makebox[12mm][l]{\fontsize{13}{16}\selectfont 54}\parbox[t]{\dimexpr\linewidth-24mm\relax}{\centering\scriptsize\bfseries DETERMINATION OF THE DIFFERENCE IN LATITUDE BETWEEN GÖTTINGEN AND ALTONA.}\makebox[12mm][r]{}\par\medskip

The observations show, however, that the graduation errors must be considerably
larger than in \textsc{Ramsden}'s zenith sector, while the observational errors
proper may even be somewhat smaller. That procedure thus gives far too much
weight to determinations based on a greater number of observations over those
based on only one or two.

If we wish to account for the influence of graduation errors, we must also consider
that each determination of the collimation error involves a constant error that
depends on the errors of the graduation marks involved. This plainly affects the
altitude of the pole in opposite directions depending on whether the circle is
east or west. We must therefore separate the observations corresponding to the
different positions of the circle, compute the mean of each series using the
weight \(\dfrac n{1+n\theta}\) for each determination, and finally take the simple
arithmetic mean of these two means.

In the absence of definite knowledge of \(\theta\), this calculation was made
under the three hypotheses \(\theta=0\), \(\theta=1\), \(\theta=\infty\), giving
the following for the altitude of the pole:
\begin{center}
\renewcommand{\arraystretch}{1.18}%
\begin{tabular}{c|c|c|c}
&\(\theta=0\)&\(\theta=1\)&\(\theta=\infty\)\\
\hline
Circle east & \(50^\circ56'5''{,}75\) & \(50^\circ56'5''{,}69\) & \(50^\circ56'5''{,}71\)\\
Circle west & \hfill \(4{,}62\) & \hfill \(4{,}65\) & \hfill \(4{,}65\)\\
\hline
Altitude of the pole & \(50\enspace56\enspace5{,}18\) & \(50\enspace56\enspace5{,}17\) & \(50\enspace56\enspace5{,}18\)
\end{tabular}
\end{center}

We thus see that applying the more rigorous principles leaves the first result
virtually unchanged, and that we can adopt the figure \(50^\circ56'5''{,}17\).

These calculations have yet to account for the flexure of the telescope. According
to Mr. \textsc{Hansen}, at the horizon it is \(=1''{,}00\), to be subtracted from
the observed zenith distance, or \(f=-1''{,}00\) in our notation. We see that
allowing for this flexure makes the altitude of the pole derived from the stars
culminating north of the zenith somewhat greater, and that derived from the
southern stars smaller. Since the former slightly out-
\par

% Printed 55; principal PDF 62; Google PDF 66.
\par\clearpage
\noindent\makebox[12mm][l]{}\makebox[\dimexpr\linewidth-24mm\relax][c]{\footnotesize\textbf{IV.\quad LATITUDE DETERMINATION FOR SEEBERG OBSERVATORY.}}\makebox[12mm][r]{\fontsize{13}{16}\selectfont 55}\par\medskip

\noindent weigh the latter, the mean result will increase by \(0''{,}02\).
The second part of the flexure, or the flexure in the vertical position, may
be regarded as a constant change in the collimation error, since all the zenith
distances occurring here are small. With our procedure it is therefore
eliminated automatically, just like the graduation errors of the marks involved
in determining the collimation error.

We accordingly have as the definitive value of the altitude of the pole from these observations
\[
50^\circ56'5''{,}19.
\]

The aforementioned trigonometric connection of the two observatories, calculated
with the dimensions of the Earth's spheroid given above, yields the difference
in latitude
\[
35'41''{,}86,
\]
\noindent and thus, together with the altitude of the pole for Göttingen determined
above, gives that of the Seeberg Observatory
\[
=50^\circ56'5''{,}99.
\]

This refers to the triangulation point, namely the centre of the axis of the
transit telescope. The centre of the axis of the meridian circle lies \(1{,}168\)
toises, or \(0''{,}07\) in arc, farther south. Thus the altitude of the pole at
the latter point, derived from Göttingen by the trigonometric connection, is
\[
=50^\circ56'5''{,}92
\]
\noindent or \(0''{,}73\) greater than that found from the astronomical observations.

The difference in longitude from the trigonometric connection is \(47'9''{,}20\)
in arc, or \(3^{\mathrm m}8^{\mathrm s}{,}61\) in time, in very good agreement
with what we know from astronomical observations. Finally, these measurements
give the azimuth of the triangle side Seeberg---southern meridian mark near
Schwabhausen as \(4''{,}6\) west. This too may be regarded as good agreement,
given the considerable number of intermediate points, the discrepancies in the
reported values of some angles in the Prussian survey, and the uncertainty
as to whether the triangulation point lay exactly on the meridian.

\begin{center}\rule{.24\textwidth}{.4pt}\end{center}

% Printed 56 (unfoliated); principal PDF 63; Google PDF 67.
\par\clearpage

\vspace*{.055\textheight}
\begin{center}{\Large ADDENDUM TO [ART. 20] P. 48.}\end{center}

\textsc{Walbeck}'s determination of the dimensions of the Earth's spheroid is
found in a short treatise, \textit{De forma et magnitudine telluris, ex dimensis
arcubus meridiani, definiendis}, of which only the first two sheets appeared in
print (Åbo, 1819). \textsc{Walbeck} included the Peruvian, the two East Indian,
the French, the English, and the more recent Lapland arc measurements in his
calculations, and to my knowledge he is so far the only person who has carried
out this task on sound principles free of arbitrariness. For each individual arc
measurement, however, he considered only the whole arc, or the altitudes of the
pole observed at the endpoints, without accounting for the intermediate points
available for several of them. In his calculation he stopped at the first power
of the flattening.

I therefore recently asked Dr. \textsc{Schmidt}, already favourably known for
several works, to undertake a new calculation of all these arc measurements.
He completed it while the last sheets of the present work were being printed.
He took into account both the higher powers of the flattening and the altitudes
of the pole observed at all intermediate points. He also included the Hanoverian
arc measurement and, following the principle indicated on p.~50 above, determined
the ellipsoid on which the astronomically observed altitudes of the pole
require the smallest possible adjustments to agree exactly with the geodetic
measurements, i.e.\ for which the sum of the squares of the required adjustments
is a minimum. The result of this calculation is:
\par

% Printed 57; principal PDF 64; Google PDF 68.
\par\clearpage
\noindent\makebox[12mm][l]{}\makebox[\dimexpr\linewidth-24mm\relax][c]{\footnotesize\textbf{DETERMINATION OF THE DIFFERENCE IN LATITUDE ETC.\quad ADDENDUM.}}\makebox[12mm][r]{\fontsize{13}{16}\selectfont 57}\par\medskip

\begin{center}
\begin{tabular}{@{}l@{\quad}l@{}}
\makebox[.62\linewidth][s]{Flattening\hfill .\hfill .\hfill .\hfill .\hfill .\hfill .\hfill .\hfill .\hfill .\hfill .\hfill .\hfill .\hfill .}
& \(\dfrac1{298{,}39}\)\\
\makebox[.62\linewidth][s]{One three-hundred-sixtieth of the Earth's meridian\hfill .}
& \(57010{,}35\) toises.
\end{tabular}
\end{center}

The following table gives the observed altitudes of the pole at the 25 points
of the seven arc measurements and the smallest adjustments they require
to agree exactly with the dimensions found for the Earth:
\[
\renewcommand{\arraystretch}{1.38}
\begin{array}{@{}l@{\quad}|r@{\enspace}r@{\enspace}r|l@{}}
\multicolumn{5}{c}{\textit{Peruvian measurement.}}\\[.3ex]
\text{Tarqui}&-3^\circ&4'&30''{,}83&+2''{,}05\\
\text{Cotchesqui}&+0&2&37{,}83&-2{,}05\\[.8ex]
\multicolumn{5}{c}{\textit{First East Indian measurement.}}\\[.3ex]
\text{Trivandeporum}&+11&44&52{,}59&-0{,}48\\
\text{Paudree}&13&19&49{,}02&+0{,}47\\[.8ex]
\multicolumn{5}{c}{\textit{Second East Indian measurement.}}\\[.3ex]
\text{Punnae}&8&9&38{,}39&-1{,}43\\
\text{Putchapolliam}&10&59&48{,}93&-1{,}18\\
\text{Dodagoontah}&12&59&59{,}91&+3{,}37\\
\text{Namthabad}&15&6&0{,}64&-0{,}77\\[.8ex]
\multicolumn{5}{c}{\textit{French measurement.}}\\[.3ex]
\text{Formentera}&38&39&56{,}11&+3{,}95\\
\text{Montjouy}&41&21&45{,}45&+2{,}81\\
\text{Barcelona}&41&22&47{,}16&+1{,}07\\
\text{Perpignan}&42&41&58{,}01&-3{,}67\\
\text{Carcassonne}&43&12&54{,}31&-0{,}96\\
\text{Evaux}&46&10&42{,}19&-6{,}14\\
\text{Pantheon}&48&50&48{,}94&-0{,}17\\
\text{Dünkirchen}&51&2&8{,}74&+3{,}12\\[.8ex]
\multicolumn{5}{c}{\textit{English measurement.}}\\[.3ex]
\text{Dunnose}&50&37&7{,}81&-1{,}73\\
\text{Greenwich}&51&28&39{,}60&+1{,}00\\
\text{Blenheim}&51&50&27{,}50&+3{,}02\\
\text{Arburyhill}&52&13&27{,}79&+1{,}80\\
\text{Clifton}&53&27&31{,}59&-4{,}07
\end{array}
\]
\noindent\textbf{IX.}\hfill 8\hspace{.12\textwidth}\par

% Printed 58; principal PDF 65; Google PDF 69.
\par\clearpage
\noindent\makebox[12mm][l]{\fontsize{13}{16}\selectfont 58}\makebox[\dimexpr\linewidth-24mm\relax][c]{\footnotesize\textbf{DETERMINATION OF THE DIFFERENCE IN LATITUDE ETC.\quad ADDENDUM.}}\makebox[12mm][r]{}\par\medskip

\[
\renewcommand{\arraystretch}{1.38}
\begin{array}{@{}l@{\quad}|r@{\enspace}r@{\enspace}r|l@{}}
\multicolumn{5}{c}{\textit{Hanoverian measurement.}}\\[.3ex]
\text{Göttingen}&51^\circ&31'&47''{,}85&-2''{,}65\\
\text{Altona}&53&32&45{,}27&+2{,}66\\[.8ex]
\multicolumn{5}{c}{\textit{Swedish measurement.}}\\[.3ex]
\text{Mallörn}&65&31&31{,}06&+1{,}40\\
\text{Pahtavara}&67&8&51{,}41&-1{,}40
\end{array}
\]

The numbers in the last column must by no means be taken as errors in the
astronomical observations. They are the algebraic sum of those errors and the
irregularities in the direction of the vertical. If we treat these total
deviations by the same rules as random errors, we find a mean deviation of
\(3''{,}18\), and hence the mean error to be expected
\begin{center}
\begin{tabular}{@{}l@{\quad}l@{}}
\makebox[.64\linewidth][s]{in the denominator of the flattening\hfill .\hfill .\hfill .\hfill .\hfill .\hfill .}
& \(12{,}5\) units\\
in the value of one three-hundred-sixtieth &\\
\makebox[.64\linewidth][r]{of the Earth's meridian} & \(5{,}0\) toises.
\end{tabular}
\end{center}

The so-called probable error may therefore be estimated at 8 units in the
denominator of the flattening and 3 toises in the mean degree of latitude.
This clarification of the degree of accuracy that may properly be attributed
to the determination of the dimensions of the Earth's spheroid from all the arc
measurements of latitude made so far is to be regarded as an important result of
this valuable work by Dr. \textsc{Schmidt}, which is to be published in detail
elsewhere.

\begin{center}\rule{.24\textwidth}{.4pt}\end{center}

% Printed 59 (unfoliated); principal PDF 66; Google PDF 70.
\par\clearpage

\vspace*{.075\textheight}
\begin{center}{\LARGE N\,O\,T\,I\,C\,E.}\end{center}
\vspace{.045\textheight}
\begin{center}
\rule{.52\textwidth}{.4pt}\\[-.2ex]
{\small Göttingische gelehrte Anzeigen.\enspace 16 June 1828.}\\[-.2ex]
\rule{.52\textwidth}{.4pt}
\end{center}

From Vandenhoeck and Ruprecht: \textit{Bestimmung des Breitenunterschiedes
zwischen den Sternwarten von Göttingen und Altona durch Beobachtungen am}
\textsc{Ramsden}\textit{schen Zenithsector, von} \textsc{Carl Friedrich Gauss}.
1828. 84 pages in 4to.

The measurements made by the author of this work in the Kingdom of Hanover
during the past few years were intended in the first instance to extend the arc measurement
undertaken by Professor \textsc{Schumacher} in the Danish states two degrees
farther south. A chain of triangles from the southern border of the Kingdom of
Hanover to Hamburg was completed in 1821---1823 and connected to the Danish
system of triangles. To derive an independent result as an arc measurement from
this chain, it was still necessary to determine the curvature of the whole
meridian arc astronomically. This, however, had to be postponed until 1827
because in the intervening years the system of triangles was extended by order
of higher authority to connect with \textsc{Krayenhoff}'s triangles. As publication
of the trigonometric work is reserved for a larger work for several reasons
that need not be set out here, Court Councillor \textsc{Gauss} took the opportunity
to issue the astronomical part of the arc measurement at once as a separate
work, so that the final result could be added immediately to the other arc
measurements and put to use.

One could not have wished for more favourably situated endpoints than those
that presented themselves here. They are the
\par\smallskip
\noindent\hfill \(8^{*}\)\hspace{.12\textwidth}\par

% Printed 60; principal PDF 67; Google PDF 71.
\par\clearpage
\noindent\makebox[12mm][l]{\fontsize{13}{16}\selectfont 60}\makebox[\dimexpr\linewidth-24mm\relax][c]{\small\textbf{NOTICE.}}\makebox[12mm][r]{}\par\medskip

\noindent observatories of Göttingen and Altona, both equipped with excellent
instruments, and both, by a singular coincidence, lying so nearly on one and
the same meridian that specific locations in the observatories must be given
to state a difference. With reference to the centres of the axes of the
\textsc{Reichenbach} meridian circles, the Altona Observatory lies only
\(7\tfrac12\) toises farther west.

The extent of the celestial arc corresponding to the portion of the Earth's
meridian between the two observatories already follows from their absolute
altitudes of the pole, for whose determination the continuing observations
constantly supply new contributions. Nevertheless, it was of great importance
to obtain a determination of the difference in latitude by another method,
with one and the same instrument of the first rank. Court Councillor \textsc{Gauss}
was able to use the excellent \textsc{Ramsden} zenith sector for this purpose.
It had been used for similar operations in the English arc measurement and,
as is well known, was described in detail by \textsc{Mudge}.

The work is divided into four sections. The first lists the stars observed.
It is characteristic of the modern art of observation that on important
occasions the number of observations is greatly increased. The aim is to
reduce the influence of the imperfections of the senses and instruments
(for none \textls[100]{can} be absolutely perfect) and of unavoidable external
disturbances to a very small fraction of its usual magnitude, since eliminating
it is impossible. This aim can be achieved only by letting the individual
sources of error come into play in many different ways over a large number of
observations. Thus, in particular, the irregular graduation errors of a
non-repeating instrument can be made to cancel out as well as possible only
by observing a considerable number of stars with different declinations.
The author therefore selected 43 stars in suitable positions for observation,
some of which do not yet appear in any star catalogue.

The second section deals with the observations themselves. They were made
in Göttingen from the beginning of April to the middle of May, and afterwards
in Altona during June, and numbered 900
\par

% Printed 61; principal PDF 68; Google PDF 72.
\par\clearpage
\noindent\makebox[12mm][l]{}\parbox[t]{\dimexpr\linewidth-24mm\relax}{\centering\scriptsize\bfseries DETERMINATION OF THE DIFFERENCE IN LATITUDE BETWEEN GÖTTINGEN AND ALTONA.}\makebox[12mm][r]{\fontsize{13}{16}\selectfont 61}\par\medskip

\noindent in all. Every observation is listed individually, arranged at once
by star together with its reduction to the mean position for the beginning of
1827. The author considered the original journal format unnecessary, and it
would have doubled the length of the work.

The third section develops the results obtained from the observations. The
principal result is the difference in latitude between the two observatories.
For the locations of the meridian circles, it is fixed at \(2^\circ0'57''{,}42\)
after several combinations from different points of view, with only a
probable uncertainty of \(0''{,}07\) remaining. The observations show most
clearly a remarkable change in the collimation error. It was \(3''{,}76\)
in the Göttingen observations and \(1''{,}39\) in the Altona observations, and
must have undergone this change during transport, although that was conducted
with the greatest care. This fact is particularly important in assessing
the confidence due to the older Lapland arc measurement. It has rightly
been faulted for failing to reverse the zenith sector at each observation
site, and the claim that the collimation error was the same is by no means
satisfactorily justified by what has recently been advanced on this point
in an otherwise very valuable treatise. One also finds here the author's
observations of 1820 and 1824, used to establish the absolute altitude of the
pole at the Göttingen Observatory, the comparison of the result for the
degree of latitude from the Hanoverian arc measurement with the dimensions
of the Earth's spheroid that \textsc{Walbeck} derived in 1819 from all reliable
earlier arc measurements, the author's views on the irregularities of the
Earth's figure generally, and finally the comparison of the declinations of
the zenith stars obtained from the present observations with the determinations
by \textsc{Bradley} and \textsc{Piazzi}. For one of the smaller stars
(\textsc{Piazzi} 16. 291), a difference of \(10''{,}3\) from \textsc{Piazzi}'s
catalogue is remarkable. For the reasons given by the author, a proper motion
is unmistakable in this difference.

The fourth section contains the determination of the latitude of Seeberg
Observatory from the observations that the director of that observatory,
\par
\par
\endgroup
